% =====================================================================
% Слайд 14 — Обучение: BC + PPO
% =====================================================================
\begin{frame}{Обучение: двустадийная схема BC $\to$ PPO}
  \footnotesize\setlength{\abovedisplayskip}{2pt}
  \setlength{\belowdisplayskip}{2pt}

  PPO с нуля плохо стартует на сложной непрерывной задаче.
  Используем стандартную схему \emph{BC $\to$ PPO $+$ KL}
  (Sequence Tutor, Jaques 2017).

  \begin{columns}[T,onlytextwidth]
    \column{0.46\textwidth}
    \textbf{Стадия A --- Behavior Cloning} на ASAP:
    \begin{enumerate}\setlength{\itemsep}{0.2em}
      \item Считаем оракульный темп
            $\displaystyle a^{*}_t = \frac{\Delta\Phi_{\text{nominal}}}{\Delta\Phi_{\text{real}}}$
            на скользящем окне.
      \item Регрессия $s_t \mapsto a^{*}_t$ --- MLP $2\!\times\!64$.
      \item Получаем стартовую политику $\pi_{\text{BC}}$.
    \end{enumerate}

    \vspace{0.4em}
    \textbf{Стадия B --- PPO $+$ KL:}
    \begin{enumerate}\setlength{\itemsep}{0.2em}
      \item Среда \texttt{ScoreFollowingEnv} (Gymnasium) поверх
            живого \texttt{HybridScoreFollower}.
      \item Симулятор солиста (rubato по Repp 1995).
      \item Награда $r_t$ из формулы\,(1).
      \item KL-штраф
            $\beta\,D_{\text{KL}}(\pi_\theta\,\|\,\pi_{\text{BC}})$
            против ухода в нефизичные регионы.
    \end{enumerate}

    \column{0.54\textwidth}
    \centering
    \resizebox{\linewidth}{!}{%
      % СХЕМА 10 — Pipeline обучения BC -> PPO
\begin{tikzpicture}[
  every node/.style={font=\scriptsize, align=center},
  box/.style={draw, thick, minimum height=0.85cm, minimum width=1.7cm,
              fill=white, rounded corners=1.5pt, inner sep=3pt},
  bigbox/.style={box, minimum height=1.1cm, minimum width=2.4cm,
                 line width=1pt, fill=black!4},
  arr/.style={-{Latex[length=2mm]}, thick},
]
  \node[box]    (asap) at (0,0)     {ASAP\\dataset};
  \node[box]    (bc)   at (2.8,0)   {BC training\\(MLP regr.)};
  \node[box]    (pbc)  at (5.6,0)   {$\pi_{\text{BC}}$\\(init)};

  \node[bigbox] (ppo)  at (5.6,-1.8) {\bfseries PPO + KL\\penalty};
  \node[box]    (env)  at (2.0,-1.8) {Score\\FollowingEnv\\+ Simulator};

  \node[box]    (pi)   at (9.0,-1.8) {$\pi_\theta$\\(final)};

  \draw[arr] (asap)  -- node[above] {$(s_t, a^{*}_t)$} (bc);
  \draw[arr] (bc)    -- node[above] {$\theta_{\text{BC}}$} (pbc);
  \draw[arr] (pbc)   -- (ppo);
  \draw[arr, latex-latex] (ppo) -- (env);
  \draw[arr] (ppo)   -- (pi);
\end{tikzpicture}
%
    }\\[0.4em]
    {\scriptsize\itshape\color{black!70}
     ASAP даёт инициализацию,\\ PPO дотачивает в симуляторе.}
  \end{columns}
\end{frame}
